%%%%%%%%%%%%%%%%%%%%%%%%%%%%%%%%%%%%%%%
% Cover Letter - НЛМК ИТ, Архитектор реинжиниринга бизнес-процессов
%%%%%%%%%%%%%%%%%%%%%%%%%%%%%%%%%%%%%%%

\documentclass[]{cover}
\usepackage{fancyhdr}

\pagestyle{fancy}
\fancyhf{}

\rfoot{Page \thepage \hspace{0pt}}
\thispagestyle{empty}
\renewcommand{\headrulewidth}{0pt}

\begin{document}

%%%%%%%%%%%%%%%%%%%%%%%%%%%%%%%%%%%%%%
%     TITLE NAME
%%%%%%%%%%%%%%%%%%%%%%%%%%%%%%%%%%%%%%
\namesection{}{\Huge{Алексей Орлов}}{  \href{mailto:lex.orlov78@yandex.ru}{lex.orlov78@yandex.ru} | +7~(901)~797-54-88 |  \urlstyle{same}\href{https://www.linkedin.com/in/alex-orlov-v/}{LinkedIn}
}

%%%%%%%%%%%%%%%%%%%%%%%%%%%%%%%%%%%%%%
%     MAIN COVER LETTER CONTENT
%%%%%%%%%%%%%%%%%%%%%%%%%%%%%%%%%%%%%%

\currentdate{\rutoday}
\lettercontent{Здравствуйте, команда НЛМК ИТ,}

\lettercontent{Меня заинтересовала позиция архитектора реинжиниринга бизнес-процессов. Из всех вакансий, что я сейчас рассматриваю, эта единственная описывает ИИ не как желательный навык, а как саму суть роли: управление полным циклом внедрения ИИ-решений от идеи до тиражирования, приоритизация по ROI, метрики успеха агентов. Ровно так я уже работаю последние годы.}

\lettercontent{Группа НЛМК ведёт цифровую трансформацию с 2019 года и в этом году планирует тиражировать подходы GenAI на все команды разработки и все этапы внедрения решений. Это именно тот момент масштабирования, где нужен не пилот на одной команде, а системный процесс, который переживёт смену конкретных инструментов.}

\lettercontent{Что я принесу на позицию:}

{\raggedright\fontspec{DejaVu Sans}\fontsize{11pt}{13pt}\selectfont
\begin{itemize}
    \item \textbf{ИИ как рабочий инструмент, не эксперимент:} использую Claude Code в ежедневном управлении процессами, а вне работы построил собственный AI-агент-шлюз (Jarvis) с MCP-интеграциями к Obsidian, Asana, Gmail и Telegram
    \item \textbf{Приоритизация по ROI:} выстроил единую модель приоритизации инициатив в ПочтаТех, сократившую Time-to-Market примерно на 20\%
    \item \textbf{Измеримые бизнес-результаты через технологические инициативы:} сократил время вывода разработок в прод примерно в 2 раза в Банке ВТБ, снизил количество конфликтов и блокеров на 35\%
    \item \textbf{Работа с владельцами процессов:} согласовывал целевую модель процесса с CTO и руководителями направлений, фасилитировал PI Planning для 15 команд / 70+ участников
\end{itemize}\par}
\vspace{6pt}

\lettercontent{Честно: формального опыта стратегического консалтинга у меня нет, весь мой опыт внедрения изменений и ИИ-инструментов был внутри компаний, не в роли внешнего консультанта. Отраслевого опыта в металлургии тоже нет. Но задача, которую вы описываете, мне знакома предметно: перевести практику из головы одного человека в процесс, который работает на масштабе всей организации. Буду рад обсудить детали.}

\begin{flushright}
\closing{С уважением,}

\signature{Алексей Орлов}
\end{flushright}
\end{document}
